% Part: many-valued-logic
% Chapter: sequent-calculus
% Section: introduction

\documentclass[../../../include/open-logic-section]{subfiles}

\begin{document}

\olfileid{mvl}{seq}{int}

\olsection{مقدمه}

حساب دنباله‌ای منطق کلاسیک دستگاه~!!{derivation} ساده و کارآمدی است.
اگر یک منطق چندارزشی با ماتریسی دارای تعداد متناهی ارزش صدق تعریف شود،
یعنی~$V$ متناهی باشد، می‌توان برای آن حسابی دنباله‌ای ارائه کرد. اندیشهٔ
این کار از بررسی معنای دنباله‌ها و شکل قواعد استنتاج در حالت کلاسیک به
دست می‌آید.

اکنون به یاد آورید که دنبالهٔ
\begin{align*}
    !A_1, \dots, !A_m & \Sequent !B_1, \dots, !B_n
\intertext{را می‌توان به‌صورت~!!{formula} زیر تفسیر کرد:}
    (!A_1 \land \cdots \land !A_m) & \lif (!B_1 \lor \cdots \lor
!B_n).
\end{align*}
به بیان دیگر، !!^a{valuation} چون~$\pAssign{v}$ دنبالهٔ
$\Gamma \Sequent \Delta$ را \emph{ارضا می‌کند} اگر و تنها اگر یا برای
یک~$!A \in \Gamma$ داشته باشیم~$\pValue{v}(!A) = \False$، یا برای
یک~$!A \in \Delta$ داشته باشیم~$\pValue{v}(!A) = \True$. بنا بر این
تفسیر، دنباله‌های آغازین~$!A \Sequent !A$ همواره ارضا می‌شوند، زیرا یا
$\pValue v(!A) = \True$ یا~$\pValue v(!A) = \False$.

قواعد استنتاج رابط شرطی در~$\Log{LK}$، با حذف فرمول‌های جانبی~$\Gamma$
و~$\Delta$، چنین‌اند:

\begin{defish}
    \Axiom$ \fCenter !A$
    \Axiom$ !B \fCenter $
    \RightLabel{\LeftR{\lif}}
    \BinaryInf$ !A \lif !B \fCenter $
    \DisplayProof
    \hfill
    \Axiom$ !A \fCenter !B$
    \RightLabel{\RightR{\lif}}
    \UnaryInf$ \fCenter !A \lif !B $
    \DisplayProof
\end{defish}

اگر تفسیر معناشناختی بالا از دنباله را به کار ببریم، قاعدهٔ
$\LeftR{\lif}$ می‌گوید اگر~$\pValue v(!A) = \True$ و
$\pValue v(!B) = \False$، آنگاه~$\pValue v(!A \lif !B) = \False$.
به‌همین‌ترتیب، قاعدهٔ~$\RightR{\lif}$ می‌گوید اگر یا
$\pValue v(!A) = \False$ یا~$\pValue v(!B) = \True$، آنگاه
$\pValue v(!A \lif !B) = \True$. در واقع، این شرطی‌ها دوشرطی‌اند. در
صورت‌بندی استاندارد قواعد~$\LeftR{\land}$ و~$\RightR{\lor}$، شرطی‌های
متناظر دوشرطی نیستند؛ اما نسخه‌های بدیلی از این قواعد وجود دارند که در
آن‌ها چنین است:

\begin{defish}
    \Axiom$!A, !B, \Gamma \fCenter \Delta$
    \RightLabel{\LeftR{\land}}
    \UnaryInf$!A \land !B, \Gamma \fCenter \Delta$
    \DisplayProof
    \hfill
    \Axiom$ \Gamma \fCenter \Delta, !A, !B$
    \RightLabel{\RightR{\lor}}
    \UnaryInf$ \Gamma \fCenter \Delta, !A \lor !B$
    \DisplayProof
\end{defish}

با اعمال این اندیشهٔ پایه به یک منطق~$n$ارزشی، حسابی دنباله‌ای با~$n$
جایگاه، به‌جای دو جایگاه، به دست می‌آید؛ یک جایگاه برای هر ارزش صدق.
برای یک منطق سه‌ارزشی با~$V = \{\False, \Undef, \True\}$، دنباله
عبارتی به‌صورت~$\Gamma \mid \Pi \mid \Delta$ است. این دنباله در
!!a{valuation} چون~$\pAssign v$ ارضا می‌شود اگر و تنها اگر یا برای یک
$!A \in \Gamma$، $\pValue{v}(!A) = \False$، یا برای یک
$!A \in \Delta$، $\pValue{v}(!A) = \True$، یا برای یک
$!A \in \Pi$، $\pValue{v}(!A) = \Undef$. در نتیجه، دنباله‌های آغازین
$!A \mid !A \mid !A$ همواره ارضا می‌شوند.

\end{document}
